\documentclass[11pt,a4paper]{article}
\usepackage[T1]{fontenc}
\usepackage[utf8]{inputenc}
\usepackage{lmodern}
\usepackage{amsmath,amssymb}
\usepackage[margin=25mm]{geometry}
\usepackage[hidelinks]{hyperref}
\hypersetup{pdftitle={A bounded Newton force, reversible preparations and a conditional action floor},pdfauthor={navstokgap research working notes}}
\newcommand{\tightlist}{\setlength{\itemsep}{0pt}\setlength{\parskip}{0pt}}
\setlength{\emergencystretch}{3em}
\setcounter{secnumdepth}{0}
\begin{document}
\section{A bounded Newton force, reversible preparations and a
conditional action
floor}\label{a-bounded-newton-force-reversible-preparations-and-a-conditional-action-floor}

\textbf{Result, 2026-10-27; written derivation, unrefereed.} A
continuous bounded Newton force can send two different exact
position--momentum states to the same state in finite time. Consequently
it admits no reversible evolution on the original point-state space. If
an admissible class of probability preparations is invariant under
affine canonical maps, closed under weak limits, and has a reversible
evolution agreeing with these incoming Newton trajectories, then its
finite-covariance preparations have a strictly positive covariance
action floor. No numerical floor is assumed in this implication.

The closure and state-covariance premises are substantive physical
assumptions. An explicit classical density-ceiling model satisfies them
and has a sharp floor. Thus the result selects a positive preparation
scale conditionally; it does not derive quantum mechanics, a universal
Planck constant, or a floor on every classical action difference.

\textbf{Targeted check (GPT-6.1 Terra, 2026-10-27).} Terra checked
sections 1--4, including singular covariances and the classical witness,
without finding a mathematical error. Its distinction between the
\(C^1\) potential and the merely continuous force is retained. The check
does not validate the physical preparation premises or a quantum
identification.

Retrieval: the \href{balaban-ode-lesson.md}{ODE lesson} constructs the
Lipschitz causal limit; the
\href{self-sufficiency-contrast.md}{self-sufficiency contrast} separates
Kepler incompleteness from exclusion of classical completions; the
\href{principia-fifth-postulate.md}{fifth-postulate note}, Theorem
B\('\) (proof read), classifies covariant Gaussian restrictions but
leaves their zero threshold available. The
\href{necessity-unit-and-indeterminacy.md}{necessity note}, Theorem I
(proof read), prevents treating an unread density ceiling as an
automatic restriction on recorded posteriors. The present argument adds
a different conditional selector: weak closure together with an actual
point-state merger. No literature priority is claimed.

\subsection{1. Existence without uniqueness, with no unbounded
force}\label{existence-without-uniqueness-with-no-unbounded-force}

Fix \(m,\gamma,R>0\), \(0<\alpha<1\), and a smooth even cutoff \(\chi\)
with \(0\le\chi\le1\), \(\chi=1\) on \([-1,1]\) and \(\chi=0\) outside
\([-2,2]\). Set

\[
 V(q)=-\frac{\gamma}{1+\alpha}|q|^{1+\alpha}\chi(q/R),
 \qquad \dot q=p/m,\quad \dot p=-V'(q).
 \tag{1}
\]

The potential is bounded and \(C^1\), and the force is continuous and
bounded. Specifically,

\[
 \|V\|_\infty\le\frac{\gamma(2R)^{1+\alpha}}{1+\alpha},\qquad
 \|V'\|_\infty\le\gamma(2R)^\alpha
 \left(1+\frac{2\|\chi'\|_\infty}{1+\alpha}\right)=:F_*.
 \tag{2}
\]

Inside \(|q|<R\), the acceleration is
\(\ddot q=(\gamma/m)\operatorname{sgn}(q)|q|^\alpha\). Its derivative
diverges at zero, so the Lipschitz hypothesis in the lesson fails.
Nevertheless solutions exist globally. Here is the relevant existence
argument without assuming uniqueness. Euler polygons on a fixed finite
time interval have uniformly bounded momentum and position, and are
equi-Lipschitz. On their common compact range the continuous vector
field is uniformly continuous, so their integral-equation residual tends
uniformly to zero with the mesh. A uniformly convergent subsequence
therefore solves (1). The bounds

\[
 |p(t)|\le |p(0)|+F_*|t|,\qquad
 |q(t)|\le |q(0)|+\frac{|p(0)|}{m}|t|+\frac{F_*t^2}{2m}
 \tag{3}
\]

prevent finite-time escape and allow successive continuation in either
time direction. This is an existence construction by compactness, rather
than the Lipschitz lesson's Cauchy construction of a unique limit.

Put

\[
 \beta=\frac{2}{1-\alpha}>2,\qquad
 B=\left[\frac{\gamma(1-\alpha)^2}{2m(1+\alpha)}
     \right]^{1/(1-\alpha)}.
 \tag{4}
\]

The curves \(q(t)=\pm B(t-t_0)^\beta\), preceded by rest at zero, solve
(1) while their positions remain inside \((-R,R)\). They join with two
continuous derivatives at \(t_0\). Rest and an arbitrary delay before
either departure are therefore all permitted by the same initial datum
\((q,p)=(0,0)\). Direct substitution uses \(\beta-2=\alpha\beta\) and
\(B^{1-\alpha}=\gamma/[m\beta(\beta-1)]\).

More decisively, choose \(s>0\) with \(Bs^\beta<R\). The two different
incoming states

\[
 z_\pm=(\pm Bs^\beta,\ \mp m\beta Bs^{\beta-1})
 \tag{5}
\]

have the unique trajectories \(q_\pm(t)=\pm B(s-t)^\beta\),
\(p_\pm=m\dot q_\pm\) for \(0\le t<s\): the vector field is locally
Lipschitz wherever \(q\ne0\). Both reach \((0,0)\) at time \(s\). These
are mergers of the full canonical state, rather than caustics of its
position projection.

\textbf{Point-state obstruction.} No injective evolution at time \(s\)
on the original \((q,p)\) states can follow both trajectories through
their arrival. A later choice of departure cannot undo the earlier
merger. An extra arrival or branch label can retain the information, but
is an enlargement of this state space.

\subsection{2. A precise preparation
theorem}\label{a-precise-preparation-theorem}

Let \(\mathcal S\) be a nonempty subset of probability measures on the
canonical plane, containing at least one measure with finite covariance.
Assume the following three properties.

\begin{enumerate}
\def\labelenumi{\arabic{enumi}.}
\tightlist
\item
  \textbf{Canonical covariance of preparations.} Every affine symplectic
  pushforward of a member of \(\mathcal S\) belongs to \(\mathcal S\).
  In one canonical pair the linear matrices have determinant one. This
  is a premise about admissible states, not a claim that the particular
  potential (1) is invariant under all these maps.
\item
  \textbf{Weak closure.} A probability-valued weak limit of a sequence
  in \(\mathcal S\) belongs to \(\mathcal S\).
\item
  \textbf{Reversible Newton compatibility.} The force (1) admits
  bijective preparation evolutions \(\mathcal E_t\) on \(\mathcal S\),
  with weakly continuous state curves. Whenever a point preparation is
  admitted, its evolution follows its unique Newton trajectory before
  arrival at the cusp. No affineness assumption on \(\mathcal E_t\) is
  needed for the conclusion below.
\end{enumerate}

For a finite-covariance preparation write

\[
 \kappa(\mu)=\sqrt{\det\operatorname{Cov}_\mu(q,p)}.
 \tag{6}
\]

This has action units. The theorem concerns covariance, including its
cross term, rather than the coordinate product alone.

\textbf{Theorem.} Under properties 1--3,

\[
 0<\kappa_*:=\inf_{\substack{\mu\in\mathcal S\\
                         \mu\text{ has finite covariance}}}
                   \kappa(\mu)<\infty.
 \tag{7}
\]

\emph{First step: point preparations are excluded.} If one Dirac measure
belongs to \(\mathcal S\), affine translations supply every Dirac
measure, in particular \(\delta_{z_+}\) and \(\delta_{z_-}\). Property 3
and weak time continuity then give

\[
 \mathcal E_s\delta_{z_+}=\delta_{(0,0)}
                        =\mathcal E_s\delta_{z_-},
 \tag{8}
\]

contradicting injectivity. Thus \(\mathcal S\) contains no Dirac
measure.

\emph{Second step: arbitrarily small covariance action would supply a
point.} For clarity fix arbitrary calibration units \(L_0,P_0>0\), with
\(S_0=L_0P_0\), and use dimensionless coordinates
\((Q,P)=(q/L_0,p/P_0)\). If (7) failed, choose preparations with
dimensionless covariance matrices \(C_n\ge0\) and
\(\sqrt{\det C_n}=\kappa(\mu_n)/S_0\to0\). Subtract each mean. Choose
\(\epsilon_n>0\) so that

\[
 \sqrt{\det D_n}\longrightarrow0,\qquad D_n=C_n+\epsilon_n I>0,
 \qquad A_n=(\det D_n)^{1/4}D_n^{-1/2}.
 \tag{9}
\]

Such a choice follows from continuity of determinant at each \(C_n\); it
also handles rank-deficient covariances. Then \(\det A_n=1\) and

\[
 A_nC_nA_n^{\mathsf T}\le\sqrt{\det D_n}\,I.
 \tag{10}
\]

These dimensionless determinant-one maps represent physical canonical
maps in the calibrated coordinates. By property 1 the centered,
transformed preparations \(\widetilde\mu_n\) remain admissible. For
every fixed \(r>0\), Chebyshev's inequality gives

\[
 \widetilde\mu_n\{Q^2+P^2>r^2\}
 \le\frac{2\sqrt{\det D_n}}{r^2}\longrightarrow0.
 \tag{11}
\]

Hence \(\widetilde\mu_n\) converges weakly to the point measure at zero.
Property 2 admits that point, contrary to the first step. The assumed
finite-covariance member also makes the infimum finite. \(\square\)

No calibrated numerical value is hidden in \(L_0,P_0\): their positive
values are arbitrary and (7) is unchanged by changing them. The state
class supplies the value of \(\kappa_*\). The theorem does not establish
that the infimum is attained, that Gaussian states are allowed, or that
all states above the floor are allowed. It needs neither a Gaussian
ansatz nor closure under adding Gaussian noise.

One may define \(\hbar_{\rm prep}=2\kappa_*>0\) as a candidate reduced
action scale. Identifying it with the phase constant \(\hbar\) needs a
separate physical argument. In quantum mechanics the familiar
Robertson--Schrödinger form is
\(\det\operatorname{Cov}(q,p)\ge\hbar^2/4\); the equality of forms alone
proves no identification. The full Planck constant would be
\(h=2\pi\hbar\), not \(\hbar\).

\subsection{3. Which assumptions do the
work?}\label{which-assumptions-do-the-work}

Without weak closure, excluding exact points leaves preparations with
arbitrarily small covariance. For example, the class of all absolutely
continuous phase-space probabilities is canonically invariant, has no
Dirac members, and contains arbitrarily narrow Gaussians. Its weak
closure includes points. Thus the conclusion is not obtained merely by
replacing points with distributions.

Without canonical squeezing, a class supported on a fixed line with a
fixed nonzero extent may be weakly closed and contain no point while
having zero covariance determinant. Canonical covariance rules out that
degenerate evasion. Without the cusp-compatibility demand, all
probabilities are weakly closed and canonically invariant and contain
points; free and harmonic dynamics transport them reversibly.

Weak closure is not a consequence of finite experimental resolution. It
is an additional completion rule admitting concentration limits whenever
they are weak probability limits. Requiring it while retaining all
arbitrarily concentrated classical preparations is exactly what would
reintroduce the forbidden point. Whether nature supplies this closure,
the full canonical preparation class, and the exact cusp force is open.
Restricting all forces to smooth Lipschitz laws removes this particular
selector.

\subsection{4. A classical model satisfies all three
premises}\label{a-classical-model-satisfies-all-three-premises}

Fix a finite density ceiling \(M>0\), of reciprocal-action units, and
let

\[
 \mathcal S_M=\{\mu:\ d\mu=\rho\,dq\,dp,\quad 0\le\rho\le M,
                                 \quad\int\rho\,dq\,dp=1\}.
 \tag{12}
\]

This class is weakly closed: if its probabilities converge weakly, their
limiting measure satisfies \(\int\phi\,d\mu\le M\int\phi\,dq\,dp\) for
every nonnegative compactly supported continuous \(\phi\), which gives
the same density bound. It is invariant under affine symplectic maps
because they preserve canonical area.

For (1), the energy-zero set is the union of the two graphs
\(p=\pm\sqrt{-2mV(q)}\) and has canonical area zero. Outside it the
Newton trajectory is unique and complete. Negative-energy paths cannot
reach zero; positive-energy paths cross it with nonzero velocity and are
uniquely continued by \(\dot q=\pm\sqrt{2(E-V(q))/m}\) near the
crossing. Away from zero the field is smooth. This full-area flow
preserves canonical area: where \(p\ne0\), energy--time coordinates give
\(dq\,dp=dE\,dt\) up to orientation, and the flow translates \(t\);
turning points away from zero have the usual smooth Hamiltonian charts.
The resulting almost-everywhere group pushes \(\mathcal S_M\)
bijectively to itself. Its state curves are weakly continuous by bounded
convergence along the continuous trajectories. No member charges the
exceptional zero-energy set, and no Dirac preparation is admitted.
Property 3 is therefore satisfied, including its stated
point-compatibility condition.

The floor can be computed sharply:

\[
 \inf_{\mu\in\mathcal S_M\text{ with finite covariance}}
             \sqrt{\det\operatorname{Cov}_\mu(q,p)}
       =\frac{1}{4\pi M}.
 \tag{13}
\]

To check it, center and canonically isotropize a positive covariance as
in section 2, now without regularization. In the calibrated plane the
density cap is \(M S_0\). Among probabilities with this cap the minimum
of \(\mathbb E(Q^2+P^2)\) is attained by filling the disk about zero at
maximum density: moving mass from a larger radius into any unfilled
smaller radius reduces the expectation. The disk has radius
\(a_0^2=1/(\pi M S_0)\) and mean square radius \(a_0^2/2\). Isotropy
gives \(2\kappa/S_0\ge1/(2\pi M S_0)\), proving the bound. A uniform
canonical ellipse, the image of this disk, attains equality. A singular
covariance would concentrate a density on a line and is impossible in
(12).

Thus the conditional theorem is consistent and genuinely weaker than
quantum kinematics. Its classical witness already carries a supplied
scale \(1/M\); the theorem says that some positive scale must appear
under its hypotheses, not how nature chooses it. It is an unread-state
model. Unrestricted classical conditioning can violate its ceiling, as
the necessity note's two-pointer construction proves. Claiming an
operational completion would also require a physical law of instruments
and records preserving its allowed states.

\subsection{\texorpdfstring{5. A supplied positive \(\hbar\) gives a
complete quantum
comparison}{5. A supplied positive \textbackslash hbar gives a complete quantum comparison}}\label{a-supplied-positive-hbar-gives-a-complete-quantum-comparison}

For every \(\hbar>0\), the same bounded real potential has

\[
 H_\hbar=-\frac{\hbar^2}{2m}\frac{d^2}{dq^2}+V(q),\qquad
 D(H_\hbar)=H^2(\mathbb R).
 \tag{14}
\]

The free operator is self-adjoint on \(H^2\). Multiplication by \(V\) is
bounded and symmetric, with relative bound zero. The Kato--Rellich
theorem therefore makes (14) self-adjoint on this same domain, with
\(H_\hbar\ge-\|V\|_\infty\). Its spectral calculus supplies the unique
strongly continuous unitary group \(U_\hbar(t)=\exp(-itH_\hbar/\hbar)\)
for every real time. These are {[}@Teschl2009, Theorem 6.4, printed
p.135, and Theorem 5.1, pp.123--124; passages read in the held first
edition{]}, with the free domain in Theorem 7.8, pp.167--168. See the
\href{../docs/Teschl_QuantumMechanics_2009.md}{source companion}.

This comparison needs neither an unbounded collision force nor a choice
of singular boundary condition. It still supplies \(\hbar>0\); the
classical model of section 4 prevents inferring that this is the only
completion. The quantum state space also is not a class of positive
joint canonical probabilities with unrestricted classical conditioning.
No such identification was used in section 2. No relativistic premise
enters, so no \(c\) is needed.

\subsection{6. Consequence for STATE}\label{consequence-for-state}

The link from Lipschitz control to Newton's polygons now has a
complement: continuous bounded forces suffice for polygon existence but
can merge full point states, obstructing reversibility there. Weak
preparation closure and canonical covariance convert that obstruction
into a conditional positive covariance action scale. The density-ceiling
control leaves quantum identification and physical record closure open.
Next test whether a record/composition principle can independently
justify or replace these closure premises, rather than assuming a
numerical floor or discarding classical enlarged-state completions.
\end{document}
